% ==========================================================================
% GENERATED by scripts/split_paper.py. Do not edit this file.
%
% Source: imported_papers/grokking/iclr2027_conference_v2.tex
% Part 6: The two ends trade robustness against adaptation speed
%
% Edits belong upstream, in the published paper. To change what the
% chapter includes, change imports/grokking.toml and re-run
% the splitter.
% ==========================================================================
\section{The two ends trade robustness against adaptation speed}
\label{sec:fragile}

Section~\ref{sec:select} leaves a rule anywhere between a concentrated circuit and a distributed one, and neither position is free. Circuits spread over many units resist damage better, and concentrated ones can be repurposed for a new task faster. We define \emph{deletion tolerance} as mean test accuracy over nine damage levels in steps of $0.1$ from $0$ to $0.8$, each zeroing the outgoing weights of a random fraction of hidden units over $30$ masks, with \emph{targeted deletion} removing the highest norm units first and \emph{frequency removal} the same construction over input embedding Fourier components (Appendix~\ref{app:repro}), random del and targeted del in tables. Under random unit deletion \citep{averbeck2006population} the distributed circuits keep much more accuracy at the same damage, at shared loss, margin and norm, on every operation and modulus we run (Figure~\ref{fig:fragile}a, Appendix~\ref{app:figs}).

Proposition~\ref{prop:fluct} says why: deleting a fixed fraction of units removes, on average, the same fraction of every circuit's summed unit contributions, so what separates them is the fluctuation around that average, which falls as $1/\sqrt{n \cdot \mathrm{PR}}$. Two deterministic attacks rule out the alternative, that random deletion misses the concentrated circuit's weak points: deleting the highest norm units first widens the gap, and a greedy search on accuracy that never looks at norms needs $23$ of $256$ units to break backpropagation's circuit against $61$ for the knob's (Appendix~\ref{app:figs}, Table~\ref{tab:plane}). The spread itself does the work: forcing the outgoing norms to a chosen participation without touching the gradient reproduces the deletion tolerance to three decimals (Appendix~\ref{app:suff}).

Adaptation speed is the other side of the trade. It is read from training time rather than from outgoing norms, so it tests the same distinction without using the measure that defined it. We grok networks on addition at each $\beta$, then fine-tune on a second operation with plain backpropagation, so only the starting circuit differs. Adaptation time rises with $\beta$ at every step, from $640$ epochs at the concentrated end to $1540$ at the distributed one, with per seed ranges that do not overlap, and only the fully distributed start fails to beat the $1410$ epochs from scratch. Part of the reason is spare capacity, backpropagation's roughly $85$ of $256$ near idle units that a new task can claim without disturbing working structure, and Appendix~\ref{app:adapt} separates that from the rest. $\beta=0.25$ matches backpropagation's source accuracy within a point while still adapting half again as slowly. Neither end is better, and choosing a rule is choosing between the two properties.

\paragraph{The effect extends beyond grokking and Fourier circuits.}
Table~\ref{tab:scope} lists every setting we applied the knob in, each removing one thing the effect could have rested on, and none matters: the operation, the group structure ($S_5$, participation from $0.70$ to $0.97$), the presence of one agreed solution the literature has characterized (sparse parity, $0.15$ to $0.96$), and the architecture (a transformer block, $0.61$ to $0.99$), with test accuracy within one percentage point of backpropagation's throughout.
